\documentclass[11pt]{article}
\usepackage[T1]{fontenc}\usepackage{newtxtext,newtxmath}
\usepackage[margin=1in]{geometry}\usepackage{parskip}
\begin{document}
\thispagestyle{empty}
\noindent Long Zhang\\
Hubei University of Automotive Technology\\
Shiyan 442002, China\\
20230100@huat.edu.cn\\[6pt]
\noindent To the Editor-in-Chief,\\
\emph{Engineering Applications of Artificial Intelligence}\\[10pt]

\noindent Dear Editor-in-Chief, \emph{Engineering Applications of Artificial Intelligence},

We submit our manuscript, ``A Data-Efficient, Differentiable Reflectarray Unit-Cell Surrogate via
Resonance-Aware Active Sampling, with a Full-Wave Anomalous-Reflection Closure,'' for consideration as a
research article.

\paragraph{In one sentence.} We introduce a \emph{resonance-aware} active-sampling methodology that learns a
sharply resonant electromagnetic response from $\sim$16 full-wave samples and supplies differentiable design
gradients---and we show it \emph{beats every off-the-shelf acquisition rule we tested}, including textbook
Gaussian-process uncertainty sampling and the standard published adaptive sampler (LOLA--Voronoi), while the
advantage \emph{transfers across three cell topologies and even to non-electromagnetic resonances}.

\paragraph{Artificial-intelligence contribution.} The core of the paper is an applied machine-learning
methodology for surrogate modelling of an expensive-to-sample, sharply resonant response. (i) We propose
a \emph{resonance-aware active-learning acquisition} that places samples by a curvature criterion and,
crucially, in a fully \emph{online} form that scores candidates using only already-acquired data---no
oracle access to the target surface. Coupled to a lightweight \emph{differentiable} neural surrogate, it
learns the full two-parameter reflection-phase response of a 24-GHz cell from about 16 samples and supplies
full-wave-consistent autodiff design gradients usable in a gradient-based inverse-design loop. We quantify
the advantage as a \emph{bootstrapped data-efficiency law}: in the resonance region the error decays with a
steeper exponent (online $N^{-1.7}$, statistically separated by non-overlapping 95\% confidence intervals
from uniform and random sampling). Critically for an applied-AI venue, the criterion is not merely better
than uniform/random: it beats both the textbook information-theoretic baseline (Gaussian-process
maximum-variance sampling, $\sim$20$\times$) and the standard published adaptive sampler for surrogate
modelling (LOLA--Voronoi, $\sim$7$\times$), because those space-fill whereas a resonant response rewards
crowding the band. The advantage also \emph{generalises}: it recurs across three reflectarray cell
topologies (patch, cross dipole, dual-resonance) and, changing only the oracle, \emph{transfers to two
non-electromagnetic resonances} (a driven-oscillator/RLC line shape and a Fano resonance)---evidence that it
is a property of the acquisition criterion, not of one device. It also \emph{widens} when a third geometry
parameter is added, provided the acquisition is stratified along the new axis---a concrete lesson on scaling
active learning to higher dimension. (ii) Through a controlled ablation we report a device-level finding:
for this resonant response a \emph{total-field} physics-informed Maxwell-residual loss is \emph{unnecessary
and counterproductive}, acting as an anti-regulariser relative to well-placed data. This holds not only
against a naive soft-constraint hybrid but against an advanced gradient-norm/NTK-balanced hybrid
($\sim$26$\times$ worse than pure data); the effective lever is data placement, not the physics residual
(we scope this to total-field formulations). These results give a transferable recipe---when active
sampling beats physics-informed regularisation for resonant surrogates---and add a quantified, device-level
case to the PINN failure-mode literature.

\paragraph{Engineering application and fit to EAAI.} The application is low-cost, gradient-based
characterisation and inverse design of reflectarray and coded-metasurface unit cells, an electromagnetic
design problem of practical importance. The pipeline is validated end to end against the openEMS full-wave
solver: an inverse-design loop hits eight target reflections to within $0.9^\circ$ on average, and an
end-to-end anomalous-reflection closure aims the reflected order to within $\approx 0.3^\circ$ of the
generalized-Snell prediction. We believe this combination---an active-learning and differentiable-surrogate
methodology, validated end to end on a real engineering task---fits the applied-AI scope of \emph{Engineering
Applications of Artificial Intelligence}.

\paragraph{Scope.} This is, by design, a simulation-only study, and we scope it honestly: the
anomalous-reflection result is reported as a proof-of-pipeline closure (correct deflection
angle; efficiency
limited by the single-resonance cell), not a finished aperture. We regard this transparent scoping as a
matter of rigor rather than a limitation of the contribution, which is methodological.

\paragraph{Assurances.} This manuscript is original, has not been published previously, and is not under
consideration by any other journal. The authors declare no competing interests. All code and data are
openly available at \texttt{https://github.com/zlera20230100/reflectarray-surrogate} (archived at Zenodo,
DOI \texttt{10.5281/zenodo.20809888}), and the manuscript includes a reproducibility table and appendices
specifying the full-wave solver configuration and the PINN residual. Highlights and a graphical abstract
accompany this submission. Generative AI was used only for language polishing; all technical content,
analysis, and conclusions are the authors' own. We are glad to suggest qualified reviewers on request.

Thank you for your consideration.

\vspace{10pt}
\noindent Sincerely,\\
Long Zhang, on behalf of the authors\\
Hubei University of Automotive Technology
\end{document}
